%%% FIELD root-test-statement
For every intervention target \(I\subseteq V\), every \(a\in I\), and every \(t\in V\setminus\{a\}\), the barred copies of \(a\) and \(t\) are marginally d-connected in \(\mathcal G_{\mathrm m,I}\) if and only if \(t\in D_I(a)\). Consequently the complete oracle \(\mathcal O_{\mathbf G}(I)\), under I-mixture faithfulness, determines \(D_I(a)\) for every \(a\in I\).
%%% FIELD root-test-proof
Fix \(a\in I\) and a copy \((a,\ell)\). By the hard-intervention construction, every component edge entering \((a,\ell)\) is deleted, and by the I-mixture-DAG construction the selector edge \(y\to(a,\ell)\) is also deleted. Thus \((a,\ell)\) has indegree zero. With the empty conditioning set, a d-connecting path has no collider. A path starting at the root \((a,\ell)\) must therefore continue along arrows directed away from \((a,\ell)\): at the first reversal, the meeting vertex would have two arrowheads and hence would be an unconditioned collider. It also cannot reach \(y\), for entering \(y\) would require traversing some edge \(y\to(v,\ell')\) backwards and would create the same first-reversal collider. Hence a marginally active path from \((a,\ell)\) to a copy of \(t\) exists exactly when \(t\) is a descendant of \(a\) in \(\mathcal G_{\ell,I}\). Taking the union over \(\ell\in L\) gives
\[
 \overline{\{a\}}\not\!\perp_{d,\mathcal G_{\mathrm m,I}}\overline{\{t\}}
 \quad\Longleftrightarrow\quad
 t\in\bigcup_{\ell\in L}\operatorname{de}_{\mathcal G_{\ell,I}}(a)\setminus\{a\}
 =D_I(a).
\]
The oracle equivalence now follows from I-mixture faithfulness and the definition of \(\mathcal O_{\mathbf G}\).
%%% FIELD full-oracle-proof
Fix \(I\subseteq V\), and abbreviate \(D(a)=D_I(a)\). Because \(\Delta=V\), every untargeted copy \((v,\ell)\), \(v\notin I\), is a child of \(y\), whereas every targeted copy is an indegree-zero root.

First let \(a,b\notin I\). For any \(\ell,\ell'\in L\),
\[
 (a,\ell)\leftarrow y\to(b,\ell')
\]
is active given \(\overline C\), since its only internal vertex \(y\) is a noncollider and \(y\notin\overline C\). Thus the barred endpoint sets are d-connected.

Next suppose \(a\in I\) and \(b\notin I\). If \(b\in D(a)\), choose a directed path from a copy of \(a\) to a copy of \(b\). If this path avoids \(\overline C\), it is active. If it meets \(\overline C\), let \(v\) be its first vertex after the copy of \(a\); then \(v\notin I\), and \(v\) has a descendant in \(\overline C\). Therefore
\[
 (a,\ell)\to v\leftarrow y\to(b,\ell')
\]
is active: \(v\) is a collider opened by its conditioned descendant. Starting instead from a path from \(a\) to some \(c\in C\), the same collider argument applies unless its first child is \(b\); in that exceptional case the one-edge prefix from \(a\) to the endpoint \(b\) is already active. Thus d-connection also holds whenever \(D(a)\cap C\ne\varnothing\). Hence
\[
 D(a)\cap(\{b\}\cup C)\ne\varnothing
 \quad\Longrightarrow\quad
 \overline{\{a\}}\not\!\perp_d\overline{\{b\}}\mid\overline C.
\]
Conversely, start from a targeted copy of \(a\) and follow any path. Its first edge, if any, points away from that copy. At the first place where the path ceases to follow a directed component path, the two incident arrowheads form a collider. Activity then requires that collider to have a descendant in \(\overline C\). Thus either the path remains directed until it reaches a copy of \(b\), giving \(b\in D(a)\), or a directed initial segment from \(a\) reaches an ancestor of a conditioned copy, giving \(D(a)\cap C\ne\varnothing\). This proves the reverse implication. The case \(a\notin I,b\in I\) is symmetric.

Finally let \(a,b\in I\). Any active path leaving the root copy of \(a\) and eventually entering the root copy of \(b\) must reverse direction at least once. The first reversal is a collider with a descendant in \(\overline C\), so \(D(a)\cap C\ne\varnothing\). Reversing the path and applying the same argument from \(b\) gives \(D(b)\cap C\ne\varnothing\). Conversely, choose component paths from copies of \(a\) and \(b\) to conditioned copies, and let \(v,w\) be their respective first vertices. If \(v\ne w\), then
\[
 (a,\ell)\to v\leftarrow y\to w\leftarrow(b,\ell')
\]
is active because both colliders have conditioned descendants. If \(v=w\), the path \((a,\ell)\to v\leftarrow(b,\ell)\) is active for the same reason. Hence
\[
 \overline{\{a\}}\not\!\perp_d\overline{\{b\}}\mid\overline C
 \quad\Longleftrightarrow\quad
 D(a)\cap C\ne\varnothing\ \text{ and }\ D(b)\cap C\ne\varnothing.
\]
D-separation of two vertex sets means separation of every pair of their vertices. Applying the three displayed criteria to all endpoint pairs therefore determines every barred-set answer, and hence the complete oracle, from \(\{D_I(v):v\in I\}\).
%%% FIELD cyclefree-proof
If \(d=0\), the union-indegree cap gives \(\operatorname{pa}_{\cup}(t)=\varnothing\) for every \(t\), so the oriented union is empty and no query is needed. Assume henceforth that \(d\ge1\).

For every ordered pair \(a\ne t\), the targeted-root descendant lemma gives
\[
 aR_{\mathbf G}t
 \quad\Longleftrightarrow\quad
 t\in D_{\{a\}}(a).
\]
Thus the \(n\) singleton queries recover \(R_{\mathbf G}\). Fix a head \(t\), and let \(A_t=\{a:aR_{\mathbf G}t\}\). Since \(\tau(\mathbf G)=0\), the directed relation \(R_{\mathbf G}\) restricted to \(A_t\) is acyclic. It therefore has an ordering in which \(u\) precedes \(a\) whenever \(aR_{\mathbf G}u\). Process candidates in this order.

Inductively let \(P_t\) be exactly the union parents of \(t\) among the candidates already processed. Before processing a new candidate \(a\), stop if \(|P_t|=d\); otherwise query
\[
 I=P_t\cup\{a\},\qquad |I|=|P_t|+1\le d.
\]
If \(a\to t\) is a true edge, it occurs in some component and survives because \(a\in I\) deletes only arrows entering \(a\), while \(t\notin I\). Hence \(t\in D_I(a)\). If \(a\to t\) is not a true edge, every component-directed path from \(a\) to \(t\) has length at least two. Its last edge is \(p\to t\) for some \(p\in\operatorname{pa}_{\cup}(t)\), and the initial part gives \(aR_{\mathbf G}p\). Thus \(p\) precedes \(a\), so the induction hypothesis gives \(p\in P_t\subseteq I\). The intervention deletes the edge entering \(p\) on that path. Every such path is cut, and therefore \(t\notin D_I(a)\). The targeted-root descendant lemma makes the test observable and proves the update invariant.

If the procedure stops at \(|P_t|=d\), the union-indegree cap excludes any further parent; otherwise every candidate is tested. Thus the final \(P_t\) equals \(\operatorname{pa}_{\cup}(t)\) for every \(t\), and the output is \(\mathsf T(\mathbf G)\). There are \(n\) singleton queries and at most \(|A_t|\le n-1\) subsequent queries for each head, hence at most
\[
 n+\sum_{t\in V}|A_t|\le n+n(n-1)
\]
paid queries.
%%% FIELD parent-upper-proof
Put \(r=\min\{d+1,n-1\}\), and let \(\mathcal F\) attain \(F(n,r,s)\). For every ordered pair \(a\ne t\), retain \(a\to t\) exactly when every \(I\in\mathcal F\) satisfying \(a\in I\) and \(t\notin I\) also satisfies \(t\in D_I(a)\); the latter event is read from the oracle by the targeted-root descendant lemma.

If \(a\to t\) is a true edge, its component copy survives every such intervention because hard intervention deletes arrows entering targets and \(t\notin I\). Therefore \(t\in D_I(a)\) in every tested row, so no true edge is discarded. Suppose instead that \(a\to t\) is not a true edge. Then \(a\notin\operatorname{pa}_{\cup}(t)\), and
\[
 |\{a\}\cup\operatorname{pa}_{\cup}(t)|\le d+1.
\]
If \(d<n-1\), extend this set inside \(V\setminus\{t\}\) to an \(r=d+1\) element set \(A\). If \(d=n-1\), absence of \(a\to t\) gives \(|\operatorname{pa}_{\cup}(t)|\le n-2\), so the same extension exists with \(r=n-1\). By the cover property there is \(I\in\mathcal F\) with \(A\subseteq I\) and \(t\notin I\). Every directed path from \(a\) to \(t\) of length at least two has a last parent \(p\in\operatorname{pa}_{\cup}(t)\subseteq I\); the edge entering \(p\) on the path is deleted. Since the length-one path is absent, \(t\notin D_I(a)\), so the decoder discards the nonedge. This proves exact recovery and
\[
 N^{\star}(n,K,d,\tau,s)\le F(n,r,s).
\]

All \(r\)-subsets form a valid cover, hence \(F(n,r,s)\le\binom nr\). For the probabilistic bound, fix \((A,t)\) with \(|A|=r\) and \(t\notin A\), and draw a uniform \(k\)-subset, where \(r\le k\le s\). Its covering probability is
\[
 p_{n,r,k}=\frac{\binom{n-r-1}{k-r}}{\binom nk}.
\]
There are \(M=n\binom{n-1}{r}\) requirements. For \(q\) independent rows, the union bound and \(1-x\le e^{-x}\) give
\[
 \Pr(\text{some requirement is uncovered})
 \le M(1-p_{n,r,k})^q
 \le M e^{-q p_{n,r,k}}<1
\]
when \(q=\left\lceil(\log M+1)/p_{n,r,k}\right\rceil\). Hence a deterministic family of at most \(q\) distinct rows exists after duplicate removal. Optimizing over \(k\), and taking the all-\(r\)-subsets construction, yields the first branch of \(U_{\mathrm{PB}}\). If \(s=n-1\), the additional \(n\) rows \(V\setminus\{t\}\), \(t\in V\), cover every requirement and yield the second branch. Therefore \(F(n,r,s)\le U_{\mathrm{PB}}(n,r,s)\).
%%% FIELD transcript-proof
Fix an exact deterministic recovery tree and follow the transcript generated by \(\mathfrak g_0\). Let \(I_1,\ldots,I_q\) be the targets on that root-to-leaf path. If, for some \(j\), none of these targets belongs to \(\mathscr D_j\), then
\[
 \mathcal O_{\mathfrak g_j}(I_h)=\mathcal O_{\mathfrak g_0}(I_h),
 \qquad h=1,\ldots,q.
\]
Induction on \(h\) shows that adaptivity chooses the same next target under \(\mathfrak g_j\) as under \(\mathfrak g_0\), so both models reach the same leaf. That leaf has one output, but \(\mathsf T(\mathfrak g_j)\ne\mathsf T(\mathfrak g_0)\), contradicting exact recovery. Consequently \(\{I_1,\ldots,I_q\}\) intersects every \(\mathscr D_j\), and \(q\) is at least the minimum hitting-set cardinality. If the nonempty families \(\mathscr D_j\) are pairwise disjoint, one target can hit at most one of them, so \(q\ge m\).
%%% FIELD faithful-proof
It is enough to prove a stronger assertion: every finite family of component DAGs on \(V\), with \(K\ge2\), has strictly positive binary parameters for which \(\Delta=V\), every component edge is active, and faithfulness holds simultaneously for all \(2^n\) interventions.

Parameterize each binary conditional kernel by its entries in \((0,1)\), each shared intervention law by \(q_i(1)\in(0,1)\), and the weights by the relative interior of the simplex. Call the resulting nonempty open parameter set \(\Theta\). Every joint atom of every \(P_I\) is a polynomial in these coordinates. For fixed disjoint \(A,B,C\), conditional independence \(X_A\perp X_B\mid X_C\) implies, for every \(a\in A\), \(b\in B\), and every binary value \(c\) of \(X_C\), the polynomial equation
\[
 p_I(0,0,c)p_I(1,1,c)-p_I(0,1,c)p_I(1,0,c)=0,
\]
where the unshown coordinates are marginalized. The cited I-mixture Markov property supplies every independence implied by d-separation. We show that each additional-independence polynomial associated with a d-connected barred pair is nonzero as a polynomial.

Because the desired parameters will have \(\Delta=V\), the complete-oracle reduction applies to the graphical criterion. A nonzero-polynomial witness may lie in the closure of \(\Theta\); nonzero evaluation there is enough to show that the polynomial is not identically zero. Give unused variables deterministic value zero and let targeted roots be independent fair bits.

If \(a,b\notin I\), put boundary weight \(1/2\) on each of two components, set \((a,b)=(0,0)\) in one and \((a,b)=(1,1)\) in the other, and set \(X_C=0\). The displayed determinant is positive. If \(a\in I,b\notin I\), choose a component-directed path from \(a\) to the first vertex \(z\in\{b\}\cup C\) reached along such a path. Its interior avoids \(C\). When \(z=b\), concentrate on that component and copy the fair bit \(a\) along the path, so \(a=b\) given \(X_C=0\). When \(z=c\in C\), use two components: in the path component copy \(a\) to \(c\) and put \(b=0\), while in the other component put \(c=0,b=1\). Given \(X_C=0\), the pairs \((0,0),(0,1),(1,1)\) have positive mass and \((1,0)\) has zero mass, again giving a nonzero determinant.

If \(a,b\in I\), choose component paths from \(a\) and \(b\) to first-hit vertices \(c,e\in C\). If the paths lie in different components, put weight on those two components, copying \(a\) to \(c\) in the first and \(b\) to \(e\) in the second while fixing the other conditioned endpoint at zero. Conditioning on \(X_C=0\) yields positive masses at \((0,1)\) and \((1,0)\) but zero mass at \((1,1)\), so the determinant is nonzero. If the paths lie in the same component and meet, copy the two root bits to their first meeting point, take their exclusive-or there, and copy it to the conditioned endpoints; conditioning on zero then makes \(a=b\). If those paths are disjoint, in that component copy \(a\) to \(c\) and \(b\) to \(e\), and mix with a second component in which \(c=e=0\). Given \(X_C=0\), this is a nontrivial mixture of the atom \((0,0)\) and an independent fair pair, whose determinant is positive. These constructions remain valid when \(c=e\): different components give the preceding two-component witness, while paths in one component meet and use the exclusive-or witness. Thus every graphically d-connected barred pair has at least one conditional-independence minor that is a nonzero polynomial.

There are finitely many choices of \(I,A,B,C\). The zero set of a nonzero real polynomial has empty interior, proved by induction on the number of variables after viewing the polynomial as a univariate polynomial with polynomial coefficients. The condition that a declared edge be inactive is also a proper polynomial equality among conditional-table entries, and the condition that a vertex fail to lie in \(\Delta\) is a proper finite family of equalities between lifted component tables. Hence the union of all extra-independence, inactive-edge, and mechanism-invariance exceptional sets has empty interior in \(\Theta\). Choose a point outside it. At that point all atom probabilities, mixture weights, and shared target probabilities are strictly positive; \(\Delta=V\); all declared edges are active; and, for every \(I\), the Markov implication together with exclusion of every extra-independence minor gives simultaneous I-mixture faithfulness. This proves the claimed realization for every lower-bound graph family.
%%% FIELD feasibility-proof
The case \(d=0\) is immediate from the union-indegree cap: every union-parent set is empty, so the target is known and \(N^{\star}(n,K,0,\tau,s)=0\).

Assume \(d\ge1\). For sufficiency when \(\tau=0\), the cycle-free policy uses targets of size at most \(d\), so it is legal whenever \(s\ge d\). For arbitrary \(\tau\), put \(r=\min\{d+1,n-1\}\). The parent-blocking proposition gives a finite nonadaptive recovery family whenever \(s\ge r\). These are exact decoders of \(\mathsf T(\mathbf G)\), establishing both sufficiency directions.

For cycle-free necessity, choose distinct vertices \(a,t\) and a set \(B\subseteq V\setminus\{a,t\}\) of size \(d-1\). In one component put the edges
\[
 a\to b\to t\qquad(b\in B),
\]
and let a second component be empty. Compare this baseline with the alternative that additionally has \(a\to t\) in the first component; repeat the empty component when \(K>2\). Their maximum union indegrees are at most \(d\), and their mixture-ancestor relations are acyclic. Give both graph families the faithful binary realizations from the realization theorem, with \(\Delta=V\). The only descendant signature that can change is that of a targeted \(a\), and
\[
 D_I^{\mathrm{base}}(a)\ne D_I^{\mathrm{alt}}(a)
 \quad\Longleftrightarrow\quad
 a\in I,\ B\subseteq I,\ t\notin I.
\]
Indeed, if some \(b\in B\) is untargeted, the path \(a\to b\to t\) masks the shortcut, while targeting \(t\) deletes the shortcut itself. The full-oracle reduction therefore makes the complete oracles equal for every \(|I|<d\), although the oriented unions differ. No finite adaptive tree using \(s<d\) can distinguish the two models.

For the cyclic necessity suppose first that \(d\le n-2\). Choose distinct \(a,b,t\) and \(C\subseteq V\setminus\{a,b,t\}\) with \(|C|=d-1\), and define
\[
 E_1=\{a\to b\}\cup\{a\to c,b\to c,c\to t:c\in C\},
\]
\[
 E_2=\{b\to a\}\cup\{a\to c,b\to c,c\to t:c\in C\}.
\]
Let the two components of \(\mathbf G\) be \(E_1\) and \(E_2\cup\{a\to t\}\), and those of \(\mathbf H\) be \(E_1\cup\{b\to t\}\) and \(E_2\); repeat an unchanged component when \(K>2\). Each component is acyclic. The union indegree of \(t\) is \(d\), that of each \(c\) is two when \(C\ne\varnothing\), and all other union indegrees are at most one; hence the cap is respected. The only directed cycle in any mixture-ancestor restriction is generated by \(aR b\) and \(bR a\), so one of \(a,b\) is a feedback vertex and \(\tau(\mathbf G)=\tau(\mathbf H)=1\).

A direct path check gives
\[
 \mathcal O_{\mathbf G}(I)\ne\mathcal O_{\mathbf H}(I)
 \quad\Longleftrightarrow\quad
 \{a,b\}\cup C\subseteq I,
 \quad t\notin I.
\]
To verify it, targeting only one of \(a,b\) leaves the opposing two-edge route through the other source; targeting both but leaving some \(c\) untargeted leaves both routes through \(c\); and targeting \(t\) deletes both disputed arrows. When all of \(\{a,b\}\cup C\) are targeted and \(t\) is not, the surviving descendant of \(a\) or \(b\) differs. The full-oracle reduction promotes these descendant statements to equality or inequality of the complete oracles. Thus every distinguishing target has size at least \(d+1\). The realization theorem makes both models legal. Hence \(N^{\star}=\infty\) when \(s\le d\) and \(1\le\tau\le n-2\).

If \(d=n-1\), the preceding acyclic pair already requires a target of size \(d=n-1\), so \(s<n-1\) is impossible even in the class with cyclic-complexity cap \(\tau\ge1\). Combining necessity and sufficiency proves the two claimed if-and-only-if statements.
%%% FIELD exact-slices-proof
We prove the three assertions separately.

For the cycle-free indegree-one class, the \(n\) singleton queries recover \(R_{\mathbf G}\) by the targeted-root descendant lemma. We claim that
\[
 a\to t\in E_{\cup}(\mathbf G)
 \quad\Longleftrightarrow\quad
 aR_{\mathbf G}t\ \text{ and there is no }k\notin\{a,t\}
 \text{ with }aR_{\mathbf G}k\text{ and }kR_{\mathbf G}t.
\]
If \(aR_{\mathbf G}t\) holds without a direct edge, a component path of length at least two supplies such a \(k\). Conversely, if \(a\to t\) is a union edge, the indegree-one cap makes \(a\) the only possible union parent of \(t\). Any path witnessing \(kR_{\mathbf G}t\) must end with \(a\to t\), and hence gives \(kR_{\mathbf G}a\). Together with \(aR_{\mathbf G}k\), this is a directed two-cycle inside the ancestor graph of \(t\), contradicting \(\tau=0\). Thus the displayed rule is exact and gives the upper bound \(n\).

For the matching lower bound, take an all-empty baseline mixture with \(\Delta=V\). For every \(a\in V\), choose \(t(a)\ne a\) so that the ordered edges \(a\to t(a)\) are distinct, and compare the baseline with a mixture having only that edge in one component. These models have union indegree at most one and cyclic complexity zero, and the realization theorem makes them legal. Under singleton targets, the full-oracle reduction shows that the alternative for \(a\) differs from the baseline exactly at \(I=\{a\}\). The \(n\) distinguishing families are nonempty and pairwise disjoint, so transcript packing gives the lower bound \(n\). Therefore \(N^{\star}_{\mathrm{cf}}(n,K,1,1)=n\).

For two directed trees on \(V=\{a,b,t\}\), use the baseline
\[
 \mathfrak g_0=(\{a\to b,a\to t\},\{b\to a,b\to t\})
\]
and the three alternatives
\[
 \mathfrak g_1=(\{a\to b,a\to t\},\{b\to a,a\to t\}),
\]
\[
 \mathfrak g_2=(\{a\to b,a\to t\},\{a\to b,b\to t\}),
\qquad
 \mathfrak g_3=(\{b\to a,a\to t\},\{b\to a,b\to t\}).
\]
Every component skeleton is a tree, and the realization theorem supplies legal faithful binary models. Among targets of size at most two, computing the targeted-root descendants and applying the full-oracle reduction gives exactly
\[
 \mathscr D_1=\{\{a,b\}\},\qquad
 \mathscr D_2=\{\{b\},\{b,t\}\},\qquad
 \mathscr D_3=\{\{a\},\{a,t\}\}.
\]
The first alternative is therefore indistinguishable under all singleton interventions, proving \(N^{\star}_{\mathrm{tree}}(3,2,1)=\infty\). When \(s=2\), the three displayed families are pairwise disjoint, so transcript packing gives a lower bound of three. Conversely query the three complements \(V\setminus\{v\}\). In the row excluding \(t\), every other vertex is targeted, so \(t\in D_{V\setminus\{t\}}(a)\) if and only if \(a\to t\) is a true edge. These three rows recover every oriented union, proving \(N^{\star}_{\mathrm{tree}}(3,2,2)=3\).

Finally consider the dense endpoint. The same complement-singleton argument gives
\[
 N^{\star}(n,K,n-1,n-2,n-1)\le n.
\]
For the reverse inequality, let the baseline have the complete DAG in the order \(1<\cdots<n\) and the complete DAG in the reverse order, repeating the latter if \(K>2\), and take \(\Delta=V\). For each \(i=2,\ldots,n\), delete \(1\to i\) from the forward component. The full-oracle reduction gives the exact distinguishing family
\[
 \mathscr D_i=\{I:\{1,\ldots,i-1\}\subseteq I,\ i\notin I\}.
\]
These \(n-1\) families are pairwise disjoint: membership in \(\mathscr D_i\) excludes \(i\), whereas membership in \(\mathscr D_j\) for \(j>i\) includes \(i\). Deleting \(n\to1\) from the reverse component gives the additional family \(\{V\setminus\{1\}\}\), disjoint from all the preceding families. All graph families have union indegree at most \(n-1\) and cyclic complexity at most \(n-2\), and the realization theorem makes them legal. Transcript packing therefore gives a lower bound of \(n\). This also covers \(n=2\), for which the ancestor restrictions have cyclic complexity zero. Hence the dense-endpoint count is exactly \(n\).
